\subsection{The exact passage from proper local sections to affine local sections}

\begin{statement}{A.20}{Proper local input and the affine conclusion}{Proved, conditional on the stated local theorem}
Fix a prime $p$ and a finite extension $k/\Q_p$. Assume the following precise input: for every smooth proper geometrically connected $k$-curve $Y$ of genus at least two, its point map
\[
Y(k)\longrightarrow\Sec(Y/k)
\]
is bijective, where sections and conjugacy use the full profinite geometric fundamental group. Then for every Fermat open curve $U_n/k$, and likewise every diagonal twist and every classical Fermat open curve of exponent at least three, every full section is point-induced or cuspidal. Over $k$ these alternatives, including the individual rational-cusp packets, are disjoint; the packet at $x$ is a torsor under $H^1(k,\wh\Z(1))$.
\end{statement}
\begin{proof}
Write $U$ for the open curve, $X$ for its smooth completion, $\Delta=\pi_1(U_{\bar k})$, and $s:G_k\to\pi_1(U)$ for a representative section. A neighbourhood is the finite etale cover corresponding to an open subgroup $H$ containing $s(G_k)$. The geometric subgroup $H\cap\Delta$ is open in $\Delta$. Since $\Delta$ is topologically finitely generated, characteristic open subgroups are cofinal among its open subgroups: intersect all subgroups of index at most the required index. There are only finitely many of these.

There is a geometric finite etale cover of $U$ whose smooth completion has genus at least two. Indeed, over $\bar k$ diagonal twists become Fermat curves, and the power cover of exponent $nm$, for $nm\ge4$, has compact genus $(nm-1)(nm-2)/2\ge3$. Choose a characteristic open subgroup $M\subset\Delta$ contained both in the subgroup of this cover and in $H\cap\Delta$. Then $Ms(G_k)$ is an open subgroup of $\pi_1(U)$ contained in $H$, projects onto $G_k$, and defines a geometrically connected neighbourhood over $k$. Its compactification has genus at least two by Riemann--Hurwitz. This constructs a cofinal system of such neighbourhoods without assuming that $s$ already lifts to the particular power cover over $k$.

For each neighbourhood $V$ in this system, project its lifted section to the proper curve $\bar V$. The assumed theorem produces a unique point $x_V\in\bar V(k)$. A morphism of neighbourhoods carries one such point to the other: the image point induces the projected section, so injectivity on the target compactification gives equality. Thus the $x_V$ form a rational pro-point in the compactified neighbourhood tower. Its image $x\in X(k)$ is independent of the level. If $x\in U(k)$, the pro-point lies in the etale tower above $x$ and its stabilizer is the point decomposition group. If $x\in X\setminus U$, its stabilizer is the decomposition group of a prolongation of that cusp. In either case $s(G_k)$ fixes this pro-point, by the neighbourhood correspondence, and hence lies in that stabilizer. Equivalently, one may apply the usual neighbourhood criterion: compatible rational points in the compactified section-neighbourhood tower give a point or a boundary decomposition subgroup. This criterion follows by taking the stabilizers of the compatible finite fibres; it uses the full inverse system, not an isolated cover.

A rational local parameter at a rational cusp splits its decomposition extension, whose geometric inertia is $\wh\Z(1)$. Changing a splitting gives a continuous $1$-cocycle, and inertia conjugation gives a coboundary. This proves the torsor description. The injection and disjointness statements follow from the argument of \citet[Theorem 14]{StixCuspidal}: pass to high-genus neighbourhoods, apply injectivity for their proper point maps, and then take the intersection of cusp stabilizers. Over a $p$-adic field proper point injectivity also follows independently from local Jacobian Kummer theory. This argument includes $p=2$ and primes dividing $n$.
\end{proof}

% The root-owned continuation A.33--A.36 is placed after the local audit
% by the integration script, preserving the Part A numerical order.

The hypothesis in A.20 is exactly Theorem 1.1 of the supplied October 6, 2026 preprint \citep{OpenAILocal2026}. We use that theorem as the requested professional-preprint input. The deduction just given verifies the affine reduction; it is not an independent proof of that preprint's completion and valuation theorems. No global consequence from its Section 8 is needed for the deduction.

\begin{statement}{A.21}{The real place and the two sides of a cusp}{Known, with the reduction proved}
For a smooth real hyperbolic curve $U$, the real section map identifies $\Sec(U/\mathbb R)$ with $\pi_0(U(\mathbb R))$. At a rational cusp, the inertia-conjugacy classes of splittings of the decomposition extension form a two-element set. Their images are the two adjacent real component classes, which can coincide when the two sides belong to one component. Different cusps can also have the same image component. Thus a disjoint packet decomposition over $\Q$ or a $p$-adic field must not be transplanted to $\mathbb R$.
\end{statement}
\begin{proof}
The section theorem for smooth quasiprojective real curves is the unnumbered introduction theorem of \citet{BrescianiVistoli2020}; these curves satisfy its hypotheses. Complex conjugation acts by $-1$ on $\wh\Z(1)$, so continuous cyclic-group cohomology gives
\[
H^1(\mathbb R,\wh\Z(1))=\wh\Z/2\wh\Z=\Z/2.
\]
A local coordinate identifies the two classes with positive and negative real tangents. Small real points with the corresponding sign lie on the adjacent real arcs, and hence induce their component classes. This also proves that the real Selmer condition adds no restriction to an actual full section. On the plus Fermat curve of even exponent there are no real points; since the affine curve is hyperbolic for $n\ge2$, its section set is empty already over $\mathbb R$.
\end{proof}

\begin{statement}{A.22}{The density gap: location, counterexample, and repair}{Proved}
The density inference reported in LM6 D.6 belongs to \citet[Corollary 7.11]{Porowski2026}, not to Section 8 of the supplied October 6 local preprint. Positive density does not imply intersection with a potentially-density-one set. Strongly positive upper density does imply that intersection. Consequently the proof of Porowski's corollary is valid with ``strongly positive'' in place of ``strictly positive''; the displayed argument does not prove its stronger printed formulation. This refutes the inference, not the geometric corollary itself.
\end{statement}
\begin{proof}
For a quadratic extension $L/K$, take $\Omega$ to be the unramified inert primes of $K$, and $\Theta$ the split primes. Chebotarev gives each density $1/2$, while $\Omega\cap\Theta=\varnothing$. The pullback of $\Theta$ to $L$ has density one: the omitted primes have residue degree two over $K$ and contribute density zero when ordered by their norms in $L$. This is precisely a potentially-density-one set disjoint from a positive-density set.

Now suppose $\Omega$ has positive upper density after pullback to every finite extension. If $\Theta$ has density one after pullback to one finite extension $L$, then its complement there has density zero and cannot contain the positive-upper-density pullback of $\Omega$. Hence $\Omega\cap\Theta$ is nonempty, and indeed has nonempty pullback after discarding any finite set.

For completeness, the corrected finiteness deduction can be made directly from the strong-density separation theorem. Suppose every $v\in\Omega$ has at most $d$ section-reduction classes, but $s_1,\dots,s_{d+1}$ are distinct Selmer sections. The finitely many pair-equality sets cover $\Omega$. If none were strongly positive, choose for each an extension making its upper density zero, and pass to their compositum. Density zero remains zero under finite extension: primes of residue degree greater than one have density zero, while degree-one primes have bounded multiplicity. The pullback of $\Omega$ would have upper density zero, a contradiction. One equality set is therefore strongly positive, and Porowski's Theorem 1.5 identifies that pair, contradicting its choice. The forward implication uses all primes. Thus the corrected bound is equivalent to the existence of such a strong-density set.

In the downloaded paper, Theorem 7.5 gives potentially density one for the comparison sets, Corollary 7.10 compares reduction cardinalities on these sets, and Corollary 7.11 makes exactly the invalid intersection inference. LM6 D.5--D.6 use the strong formulation. The Fermat proof uses finite descent and Stoll's density-one theorem instead; it does not use this positive-density corollary. The supplied local preprint's Section 8 does not state the disputed corollary.
\end{proof}

\subsection{Twist classes and the correct meaning of a finite answer}

\begin{statement}{A.23}{Exact Kummer fibres of cusp packets}{Proved}
Let $B=\mathbb P^1_\Q\setminus\{0,1,\infty\}$, let $n\ge3$, and let $[\alpha],[\beta]\in\Q^\times/\Q^{\times n}$. The corresponding diagonal cover is
\[
V_{\alpha,\beta,n}:\quad \alpha X^n+\beta Y^n=Z^n,\qquad XYZ\ne0,
\quad t=\alpha(X/Z)^n.
\]
Write $r_n:\wh{\Q^\times}\to\Q^\times/\Q^{\times n}$ for reduction. Using the parameters $t,1-t,1/t$ at $0,1,\infty$, respectively, its Kummer fibre on the tripod is exactly
\[
\begin{array}{c|c|c}
\text{cusp}&\text{condition on the classes}&\text{parameters in its packet}\\\hline
0&[\beta]=1&r_n(c)=[\alpha]\\
1&[\alpha]=1&r_n(c)=[\beta]\\
\infty&[\beta]=[-\alpha]&r_n(c)=[\alpha]^{-1}.
\end{array}
\]
Every nonempty parameter set in this table is infinite, indeed uncountable. Therefore the additional assertion in task A.5 that there are finitely many section classes is false, already for $\alpha=\beta=1$ and $n=3$.
\end{statement}
\begin{proof}
Over $B$, adjoining $n$th roots of $t/\alpha$ and $(1-t)/\beta$ is a finite etale torsor under $\mu_n^2$. Pulling this torsor back along a section trivializes it exactly when the two residual Kummer classes equal $[\alpha],[\beta]$. A chosen trivialization gives a lift to the displayed twist. At a cusp the unit factors in $1-t$, $t$, or $1-1/t$ have value one, so their tangential Kummer classes vanish. The pairs of completed Kummer classes are consequently
\[
(c,1),\qquad (1,c),\qquad(c^{-1},-c^{-1}).
\]
Reduction modulo $n$ proves every entry. Geometric points above these cusps satisfy respectively $\beta Y^n=Z^n$, $\alpha X^n=Z^n$, and $\alpha X^n=-\beta Y^n$, giving the same rationality conditions. For even $n$ a nonzero rational $n$th root has two signs, so counting rational cusps upstairs and counting downstairs packets are different operations.

To prove infinitude without any topological cardinality inference, take a prime $r$ and the distinct rational Kummer parameters $c_j=c_0r^{nj}$ for integers $j$. Their valuations at $r$ distinguish their classes in $\wh{\Q^\times}$; their reductions modulo $n$ agree. Allowing $j\in\wh\Z$ gives an injection of $\wh\Z$ into this fibre. The packet injection over $\Q$ is Stix's Theorem 14. All these sections are Selmer because they are globally cuspidal. For the base cubic fibre, the main Fermat theorem shows that the three listed packets, with parameters reducing to one, are the complete answer; this is an infinite union of three infinite fibres, not a finite set of sections.
\end{proof}

\begin{statement}{A.24}{Rank-zero target criterion for every diagonal twist}{Proved under the displayed arithmetic condition}
Let $V/\Q$ be a hyperbolic open curve with smooth completion $X$. Suppose that $X$ admits a nonconstant morphism to a genus-one curve $T/\Q$ whose Jacobian $E$ satisfies
\[
\#E(\Q)<\infty,\qquad \Sha(E/\Q)_{\mathrm{div}}=0.
\tag{A24.1}
\]
Then
\[
\Sec^{\mathrm{Sel}}(V/\Q)
=\{s_x:x\in V(\Q)\}\ \sqcup\!
\coprod_{c\in(X\setminus V)(\Q)}\mathcal P_c,
\tag{A24.2}
\]
where $V(\Q)$ and the indexing set of cusps are finite, and $\mathcal P_c$ is the full infinite cusp packet. A rational point on $T$ is not required. With A.20's local theorem, the same equality holds for all full sections.
\end{statement}
\begin{proof}
Choose a rational divisor $D$ of positive degree $d$ on $T$, for instance a hyperplane section of a projective embedding. The rational morphism
\[
T\longrightarrow E=\operatorname{Pic}^0(T),\qquad x\longmapsto[d x-D]
\]
is defined over $\Q$. After choosing a geometric point on $T$ it is multiplication by $d$ followed by translation, so it is a finite nonconstant morphism of degree $d^2$. Thus $X$ maps nonconstantly to $E$. Stoll's Corollary 8.1 identifies the finite-abelian-descent locus on $E$ with the closure of $E(\Q)$ under (A24.1); this closure equals $E(\Q)$ because the group is finite. Full finite descent is contained in that locus and contains the rational points. The finite-support lemma in the main audit therefore applies to $X\to E$. It gives precisely the forward inclusion in (A24.2); rational point and cusp sections give the reverse inclusion. The finite morphism $X\to E$ has finite fibres, so $X(\Q)$ is finite. Section injectivity and cusp disjointness over $\Q$ identify the indicated union without multiplicities.
\end{proof}

For a plane diagonal cubic, the use of a degree-three divisor in this proof avoids an unjustified identification of a genus-one torsor with its Jacobian before a rational point has been found. For any elliptic curve over $\Q$ with complex multiplication, the additional condition $L(E,1)\ne0$ implies (A24.1): \citet[Corollary C and Remark (3), p.~528]{AMRubin1987} give rank zero and finiteness of $\Sha(E/\Q)$, respectively. We opened these passages in the original article; Theorem A, p.~527, is first stated over the imaginary quadratic CM field, whereas Remark (3) gives exactly the required conclusion over $\Q$. Finite generation then turns rank zero into finiteness of $E(\Q)$. This proves Lead 2 for the cubic, Gaussian and sextic targets displayed here. It does not compute a missing $L$-value, and the word ``finite'' applies to rational points and cusp locations, not to cusp section classes.

\begin{statement}{A.25}{The three quartic quotient targets on an arbitrary twist}{Proved}
For $a,b,c\in\Q^\times$ and $X:a x^4+b y^4=c z^4$, the three quotient maps named in task A.4 have equations
\[
\begin{array}{c|c|c}
(u,v)&\text{genus-one target}&\text{Jacobian }Y^2=X^3+AX\\\hline
(x/y,z^2/y^2)&v^2=(a/c)u^4+b/c&A=-4ab/c^2\\
(x/z,y^2/z^2)&v^2=-(a/b)u^4+c/b&A=4ac/b^2\\
(y/z,x^2/z^2)&v^2=-(b/a)u^4+c/a&A=4bc/a^2.
\end{array}
\tag{A25.1}
\]
All geometric endomorphism rings of these Jacobians are $\Z[i]$. If any one of the three Jacobians satisfies (A24.1), the Selmer section set of $X\setminus\{xyz=0\}$ is (A24.2). Requiring all three to have analytic rank zero is a sufficient, stronger condition once the rank-zero theorem has been verified; there is no $\Q(\sqrt{-2})$ factor in these three targets.
\end{statement}
\begin{proof}
Substitution gives the equations and the maps extend to finite maps between smooth projective curves. For $v^2=A_0u^4+B_0$, choose $b_0^2=B_0$. The birational map with origin $(0,b_0)$ is
\[
X=\frac{2b_0(v+b_0)}{u^2},\qquad Y=\frac{4b_0^2(v+b_0)}{u^3}.
\]
Direct substitution gives $Y^2=X^3-4A_0B_0X$, and $dX/(2Y)=-du/(2v)$, independent of the sign of $b_0$. The induced Jacobian isomorphisms therefore have Galois-invariant differential and preserve the identity, so they descend. This proves the formula over $\Q$ even when the quartic has no rational point. Its discriminant is nonzero, and $j=1728$. The automorphism $(X,Y)\mapsto(-X,iY)$ has square $[-1]$, so the geometric endomorphism ring contains $\Z[i]$. In characteristic zero the geometric endomorphism algebra is either $\Q$ or an imaginary quadratic field; since $\Z[i]$ is already the maximal order of $\Q(i)$, the ring is exactly $\Z[i]$. Apply A.24 to any one quotient. This gives a criterion, not the exhaustive analytic-rank census requested in A.7.
\end{proof}

\begin{statement}{A.26}{Explicit sextic quotient criteria}{Proved}
For $V_{\alpha,\beta,6}$ there is a finite etale map to $V_{\alpha,\beta,3}$ given by $(X,Y,Z)\mapsto(X^2,Y^2,Z^2)$, extending to a nonconstant map of completions. In addition its completion maps nonconstantly to the elliptic curves
\[
E_1:y^2=x^3+\alpha^2\beta^3,
\qquad E_2:y^2=x^3+\alpha^3\beta^2.
\tag{A26.1}
\]
Any one of these elliptic targets, or the Jacobian of the cubic target, satisfying (A24.1) forces every Selmer section of the sextic twist to be point-induced or cuspidal, with the complete geometric description (A24.2).
\end{statement}
\begin{proof}
On $Z=1$ put $u=X^2$, $v=Y^3$. Then $\alpha u^3+\beta v^2=1$. The functions
\[
x=-\alpha\beta u,\qquad y=\alpha\beta^2v
\]
satisfy the first equation of (A26.1). Interchanging $X$ and $Y$ gives the second. The functions are nonconstant, and properness extends each map across the omitted points. The cubic map is etale on the open curve because its only possible power-map ramification lies on the coordinate axes. All maps preserve local point or cusp sections after compactification, so A.24 applies. These are sufficient quotient criteria, not a claim that every elliptic quotient has been enumerated.
\end{proof}

\begin{statement}{A.27}{Jacobian of a diagonal cubic}{Proved}
For nonzero $a,b,c\in\Q$, the Jacobian of $aX^3+bY^3+cZ^3=0$ is
\[
E_{abc}:\quad y^2=x^3-432(abc)^2,
\]
equivalently the elliptic curve $u^3+v^3=abc$ with its rational point at infinity. Its isomorphism class depends on the cube class of $abc$. In the normalization of A.23, the sign disappears, and the rank-zero criterion depends on $[\alpha\beta]\in\Q^\times/\Q^{\times3}$.
\end{statement}
\begin{proof}
Choose cube roots $A^3=a$, $B^3=b$, $C^3=c$ and put $\delta=ABC$. Scaling the projective coordinates by $A,B,C$ identifies the curve with the Fermat cubic over the splitting field. On $Z=1$ the regular differential $\omega=dx/(3by^2)$ is carried to the Fermat differential with factor $AC/B^2=\delta/b$. Equivalently, the pullback of $dx'/(3y'^2)$ is $\delta\omega$. The Fermat Jacobian is $y^2=x^3-432$, as follows from the explicit birational map in the main audit. Scaling this Weierstrass model by $(x,y)\mapsto(\delta^2x,\delta^3y)$ produces $E_{abc}$ and scales its invariant differential by $\delta^{-1}$. Hence the composite induced isomorphism of Jacobians preserves a fixed rational differential, up to a constant independent of the chosen cube roots. Its Galois conjugates have the same differential and preserve the identity. An automorphism of an elliptic curve in characteristic zero with identity differential is the identity, so this isomorphism descends to $\Q$. Finally changing $abc$ by $r^3$ gives the rational scaling $(x,y)\mapsto(r^2x,r^3y)$; replacing it by its negative has no effect on the displayed Weierstrass equation.
\end{proof}

\subsection{What prime-to-residue-characteristic ramification actually proves}

\begin{statement}{A.28}{Simultaneous cusp recognition at good primes}{Proved}
Let $n\ge3$ be odd, let $q\nmid n$ be a rational prime, and let $s$ be any full section of $U_n/\Q$, without a Selmer hypothesis. In the local geometric maximal prime-to-$q$ quotient, if the inertia image of $s|_{G_{\Q_q}}$ is nontrivial, the section lies in the decomposition group of one cusp in $(X_n\setminus U_n)(\Q_q)$. In particular, if its geometric pro-$\ell$ projection is ramified for any $\ell\ne q$, all its pro-$\ell'$ projections, for $\ell'\ne q$, are images of the same local cusp section. This proves the requested prime-to-$q$ assertion at every good prime simultaneously; it is stronger than separate pro-$\ell$ classifications but does not treat $q\mid n$.
\end{statement}
\begin{proof}
Use an integral tangent at a rational cusp to fix a splitting and genuine arithmetic action. The projective Fermat model is smooth over $\Z_q$ because $q\nmid n$, and its boundary is finite etale and disjoint. Prime-to-$q$ specialization identifies the geometric prime-to-$q$ fundamental group $P$ with that of the punctured special fibre, with unramified genuine action. Its group is free in the category of pro-groups of finite order prime to $q$; this category is closed under subgroups, quotients and extensions. The boundary has at least three geometric points. Every boundary inertia subgroup is a procyclic free factor, and any two distinct cusp generators occur together in a free basis after eliminating a third boundary generator.

Theorems A' and B' of \citet{cHerfortRibes1985}, applied to this extension-closed category, show that these factors are torsion-free, malnormal and equal to the centralizer of each of their nonidentity elements. Distinct cusp factors have nontrivial intersection with no conjugates of one another: their images are independent in the free abelianization. Consequently any abelian subgroup meeting a cusp factor nontrivially lies in it. The union $W$ of all conjugates of all cusp factors is compact, being a finite union of images of compact sets.

Wild inertia acts trivially on $P$ and any homomorphism from its pro-$q$ group to $P$ is trivial. The remaining image of inertia under the section is procyclic, say $I$, and the section's Frobenius $\psi$ acts on it by $q$th powers. Suppose $I\ne1$ and it lies in no cusp factor. Choose $\ell\ne q$ for which $I_\ell\simeq\Z_\ell$ is nonzero. The preceding centralizer assertion gives $I_\ell\cap W=\{1\}$. The compact annulus $I_\ell\setminus\ell I_\ell$ and $W$ can be separated in a finite quotient of $P$. Shrinking its kernel to a characteristic open subgroup $N$ preserves this separation. In particular $I_\ell\cap N\subset\ell I_\ell$.

The open subgroup $H=NI_\ell$ is $\psi$-stable. The quotient character
\[
H\longrightarrow I_\ell/(I_\ell\cap N)
\longrightarrow I_\ell/\ell I_\ell\simeq\F_\ell
\]
is nonzero on $I_\ell$ and kills every boundary inertia subgroup of $H$, by the separation. The arithmetic subgroup generated by $H$ and $\psi$ defines a geometrically connected finite etale cover of the punctured special fibre. Here the continuous map $\wh\Z\to P\rtimes\wh\Z$ sending $1$ to $\psi$ is a splitting, so this descent is over the actual residue field. Let $Y$ be the smooth compactification of that cover. Filling its boundary kills exactly those inertia groups, and its geometric pro-$\ell$ abelianization is $T_\ell\operatorname{Jac}(Y)$. Thus the character detects a nonzero image of $I_\ell$ in that torsion-free Tate module. Frobenius acts on that image by $q$. This contradicts the finite-field weight theorem, whose eigenvalues have absolute value $\sqrt q$; see \citet[II, Theorem 1.1(b); III, Proposition 9.1]{Milne2008}.

Therefore $I$ lies in a unique cusp inertia subgroup. Since the whole local section normalizes $I$, malnormality puts it in the corresponding cusp decomposition group. Surjectivity onto $G_{\Q_q}$ implies that the cusp is $\Q_q$-rational. Projection to any pro-$\ell'$ quotient gives the asserted simultaneous conclusion. The proof uses a finite geometric cover detecting the entire inertia relation; merely checking the relation in abelianization would not suffice.
\end{proof}

\begin{statement}{A.29}{The precise cross-prime parameter statement}{Proved; an integral interpretation is refuted}
Under A.28, the common local cusp carries a single prime-to-$q$ completed Kummer parameter, whose $\ell$-components give the individual pro-$\ell$ parameters. This does not mean that the valuation parameters for different $\ell$ are projections of one ordinary integer or that ramification occurs in every pro-$\ell$ quotient. Both stronger interpretations of Lead 4 are false, even for globally cuspidal sections.
\end{statement}
\begin{proof}
For $\ell\ne q$, local Kummer theory gives, after fixing a uniformizer and a tame generator,
\[
H^1(\Q_q,\Z_\ell(1))\simeq
\Z_\ell\oplus\Z_\ell/(q-1)\Z_\ell.
\]
The first coordinate is the valuation; the second comes from the $\ell$-part of the finite residue-field unit group. The pro-$q$ principal units disappear. Taking the product over $\ell\ne q$ gives the prime-to-$q$ completion of $\Q_q^\times$. Thus compatibility is exactly projection from this completed group, with no requirement of diagonality in an ordinary integer.

For example, take $q=7$ and let $e\in\wh\Z$ have $e_3=1$, $e_5=2$, and $e_\ell=0$ for every other prime. Twist a rational-cusp splitting of any $U_n/\Q$ by the global completed Kummer class $7^e$. The valuation at $7$ in its pro-$3$ projection is $1$, in its pro-$5$ projection it is $2$, and in its pro-$2$ projection it is zero. No integer has these components, since an integer with zero $2$-adic component is zero. Nevertheless the class exists in $\varprojlim_m\Q^\times/\Q^{\times m}$: use compatible residue classes of $e$ to define $7^{e\bmod m}$. It gives an actual full cuspidal section, not a collection of unrelated local cocycles. All nonzero projections have the same cusp. This separates the valid common-cusp conclusion from the invalid common-integer conclusion.
\end{proof}

\begin{statement}{A.30}{The remaining local gap and a proof-level sufficient hypothesis}{Open in the unrestricted form requested}
A.28 does not independently prove task A.9 at primes dividing $n$, and does not force an unramified local section to be point-induced. It also does not identify the $q$-primary quotient at $q$. Consequently the proposed weakening of the local input in task A.11 and Lead 4, with all bad primes included, is not established by prime-to-$q$ specialization. The literal conclusions of tasks A.9--A.10 do follow using the stipulated local theorem, as A.31 records.
A sufficient hypothesis for the modular Fermat proof is the conjunction of: (i) the odd residual Legendre representation admits a quadratic twist unramified outside $2\ell$, finite flat at $\ell$, of Serre conductor at $2$ at most $4$; (ii) after choosing its invariant line, the resulting section on $Z_\ell$ and its neighbourhoods supply an adelic localization surviving finite descent, with the finite-support condition imposed only at density-one finite places. The full local theorem in A.20 supplies both through the corrected local calculations. No assertion is made that this conjunction is logically minimal among all possible proofs.
\end{statement}
\noindent\textit{Exact obstruction and attempts.}
At $q\mid n$ the model used in A.28 is not smooth with etale boundary, so the displayed comparison with a smooth finite-field Jacobian is unavailable. Passing to semistable reduction introduces components and nodes; a weight-two inertia direction need not come from an original puncture. Applying the tripod comparison does not repair the issue for every coefficient prime: geometric pro-$\ell$ completion need not preserve injection for covers whose degree is not an $\ell$-power. At unramified good primes the hypothesis $I\ne1$ is absent, so the compact-separation argument has no nonzero character to detect. A.20 resolves these issues conditional on its proper local input, but an independent deduction of the ``weakest'' hypothesis from only the proposed prime-by-prime data has not been obtained.

\begin{statement}{A.31}{All-prime compatibility using the stipulated local theorem}{Proved with the named local input}
Assume Theorem 1.1 of \citet{OpenAILocal2026} in the form of A.20. Every full section of the odd plus Fermat curve $U_n/\Q$, $n\ge3$, is globally cuspidal. Consequently, at every prime $q$, including $q\mid n$, its image in every geometrically pro-$\ell$ quotient is the image of one and the same full local cusp section. After choosing a rational tangent at its global cusp, all these parameters are projections of a single element of $\wh{\Q^\times}$, localized in $\wh{\Q_q^\times}$. In particular the ramified-case assertion of task A.9 holds, and task A.10 has the completed-parameter meaning of A.29. The same conclusion holds for every classical exponent $n\ge3$.
\end{statement}
\begin{proof}
At every finite place A.20 makes the localization point-induced or cuspidal; A.21 supplies the real condition. The section is therefore Selmer in the precise sense of A.1. The repaired odd-exponent theorem and its composite reduction show that it belongs to a rational global cusp packet. For the classical even exponents we use the quartic and power-cover reductions of A.16. A rational tangent splits the global cusp decomposition extension with geometric kernel $\wh\Z(1)$. Thus a representative of the section is obtained by a continuous $1$-cocycle in $H^1(\Q,\wh\Z(1))=\varprojlim_m\Q^\times/\Q^{\times m}$. Restriction to a decomposition group and projection of coefficients commute. They give the single localized parameter and every coefficient-prime projection asserted. This uses no good-reduction assumption and makes no assertion that the completed value is an ordinary rational number or integer. It uses the local theorem to obtain global cuspidality, and hence cannot be used as an independent replacement for that theorem in the proof of global cuspidality.
\end{proof}

\begin{statement}{A.32}{The exact consequences for LM6 B.27 and AN D.30}{Proved, with the scope stated}
Let a full section of the rational tripod have a cubic, quartic or sextic twist class satisfying one of A.24--A.27's rank-zero criteria. Under the local input of A.20 it is point-induced or cuspidal. If its geometric cyclotomic centralizer is trivial, it is point-induced. Its pulled-back Legendre Tate system is then ramified at only finitely many rational primes. Without the local theorem the same statements hold for Selmer sections. These conclusions settle the indicated twist classes over $\Q$ in LM6 B.27 and AN D.30; they do not settle arbitrary classes, arbitrary number fields, or the finite-ramification question for all noncuspidal sections.
\end{statement}
\begin{proof}
The prescribed Kummer class lifts the section to the corresponding finite etale diagonal cover. Selmer localization is preserved by this lift: the local section-neighbourhood criterion, or direct restriction of a point or cusp decomposition group, gives a point or cusp upstairs whenever the lifting section exists. For all sections, apply A.20 to that cover instead. The finite-support criterion A.24 then gives a point or cusp upstairs, which projects to a point or cusp downstairs.

The cyclotomic centralizer here is the subgroup of the geometric fundamental group commuting with $s(\ker\chi)$. A cuspidal section acts on its nonzero geometric inertia $\wh\Z(1)$ by $\chi$, even after changing its cusp cocycle, since the inertia group is abelian. Thus this entire inertia group lies in the cyclotomic centralizer. Triviality excludes the cusp alternative. For a rational point $t\in\Q\setminus\{0,1\}$ the pulled-back family is the actual elliptic curve $y^2=x(x-1)(x-t)$. An integral model has nonzero discriminant, so outside the finitely many primes dividing its denominators and discriminant it is smooth; its prime-to-residue-characteristic Tate modules are unramified by smooth proper base change. This is the required finite-ramification assertion.

There is no analogous assertion for all cuspidal sections. To see the issue with an actual full section, enumerate the primes other than $2$ and put
\[
c=\prod_{j\ge1}p_j^{j!}\in\varprojlim_m\Q^\times/\Q^{\times m}.
\]
For each $m$, all but finitely many exponents vanish modulo $m$, so this compatible element is well defined. Use it as the tangent parameter at $0$ of the tripod. For any fixed coefficient prime $\ell\ge5$, its Legendre extension class is $2\kappa_\ell(c)-8\kappa_\ell(2)$. At every $p_j\ne\ell$ its inertia value is $2j!\ne0$ in $\Z_\ell$. Thus the rational Tate representation is ramified at infinitely many primes, although every finite residual representation is ramified at only finitely many. This example is cuspidal and does not refute AN D.30's explicitly noncuspidal statement. It proves why that qualifier cannot be dropped. We do not infer the outstanding general statements from a union of rank-zero twist classes.
\end{proof}
